\section{床沈み込みが歩行制御に与える影響}

床沈み込みは，二足歩行ロボットの歩行制御に対して複数の影響を与える．第一に，支持脚足先の実際の接地位置が，計画された位置からずれることである．この位置誤差は，運動学的拘束条件の破綻を引き起こし，関節角や姿勢制御に影響を及ぼす．

第二に，床反力の立ち上がり特性が変化する点が挙げられる．剛体床では，接地と同時に十分な床反力が得られるのに対し，軟弱地面では床変形を伴うため，床反力の立ち上がりが遅れ，支持力が一時的に不足する可能性がある．この影響は，特に片脚支持期において顕著であり，転倒リスクの増大につながる．

第三に，接地タイミングのずれや床反力分布の変化により，想定されたZMP軌道からの逸脱が生じる．これらの影響は，従来の歩行生成・安定化制御手法においては十分に補償されない場合が多い．